\documentclass[letterpaper]{article}
\usepackage[utf8]{inputenc}
\usepackage[margin=1in]{geometry}
\usepackage{amsmath,amssymb,amsfonts}
\usepackage{booktabs}
\usepackage{multirow}
\usepackage{graphicx}
\usepackage{hyperref}
\usepackage{natbib}
\usepackage{xcolor}
\usepackage{enumitem}
\usepackage{caption}
\usepackage{subcaption}
\usepackage{array}
\usepackage{tabularx}

\title{ContamCal: Calibrating LLM Benchmark Scores\\via Contamination-Augmented Item Response Theory}

\author{
Anonymous Authors
}

\date{}

\begin{document}
\maketitle

%==============================================================================
\begin{abstract}
%==============================================================================
Benchmark contamination---where test data leaks into training corpora---inflates LLM evaluation scores by up to 11--15 percentage points and renders model comparisons unreliable.
Existing mitigation strategies (rephrasing, option shuffling) are systematically ineffective, while detection methods flag contamination without quantifying its effect on scores.
We introduce \textbf{ContamCal}, a two-stage framework that (1)~estimates per-item contamination exposure via reciprocal rank fusion of four detection signals, and (2)~recovers calibrated benchmark scores with Bayesian credible intervals using a contamination-augmented 4PL IRT model with a measurement error component for noisy exposure estimates.
On controlled SFT contamination injection experiments across Qwen2.5 (0.5B--7B) and Llama-3.2-1B, ContamCal achieves 3.1$\pm$0.5\,pp MAE with 87\% coverage at 95\% nominal level under the SFT regime, significantly outperforming seven baselines including DCR (4.1\,pp, $p{=}0.042$) and remove-and-rescore (5.6\,pp, $p{=}0.008$).
Under RL post-training, where detection signals degrade severely (RRF AUC: $0.76 \to 0.56$), ContamCal produces appropriately wide credible intervals reflecting genuine epistemic uncertainty, reducing expected selection regret by 55\%.
A systematic signal$\times$regime AUC analysis reveals that signal survival is regime-conditional: embedding and Min-K\%++ signals collapse post-RL while Self-Critique partially survives---an empirically novel finding with direct implications for the contamination detection community.
Fisher Information-grounded identifiability analysis establishes AUC${\geq}0.65$ as the practical threshold for reliable estimation, providing a diagnostic tool for practitioners to assess when calibration results are trustworthy.
\end{abstract}

%==============================================================================
\section{Introduction}
\label{sec:intro}
%==============================================================================

Benchmark-driven evaluation is foundational to LLM development, yet increasingly compromised by training data contamination.
When test items leak into training corpora---whether through direct inclusion, paraphrased variants, or data pipeline accidents---benchmark scores inflate substantially, with controlled experiments showing 8--15 percentage point gains on MMLU alone~\citep{yang2023rethinking,oren2024proving}.
The scale of this problem is no longer hypothetical: community-wide analyses suggest contamination affects the majority of popular benchmarks~\citep{deng2024investigating,balloccu2024leak}.

The natural response---mitigating contamination through test set modification---has proven systematically ineffective.
\citet{xu2024benchmarking} tested 20 distinct mitigation strategies (rephrasing, option shuffling, difficulty balancing) and found that none reliably eliminates the inflation signal.
This ``Emperor's New Clothes'' finding motivates a fundamentally different approach: rather than attempting to eliminate contamination (which does not work), we statistically control for it---measuring the contamination effect and subtracting it with proper uncertainty quantification.
This parallels the shift in causal inference from eliminating confounders to statistically adjusting for them.

We introduce \textbf{ContamCal}, a two-stage framework grounded in educational measurement's item compromise IRT tradition~\citep{eckerly2017detecting,sinharay2017assessing} and explanatory IRT~\citep{deboeck2004explanatory}.
Stage~1 computes per-item contamination exposure scores by fusing four detection signals (embedding retrieval, Min-K\%++, ConStat, Self-Critique) via reciprocal rank fusion (RRF)~\citep{cormack2009reciprocal}---a deliberate design choice that avoids meta-circularity risks inherent in learned aggregation.
Stage~2 feeds these exposure estimates into a contamination-augmented 4PL IRT model with per-benchmark contamination coefficients $\gamma_k$ and a measurement error model that propagates detection noise into posterior uncertainty.
The output is a calibrated benchmark score with Bayesian credible intervals that are regime-conditional: tight under SFT (where signals are reliable) and appropriately wide under RL (where signals degrade).

Our contributions are:
\begin{enumerate}[leftmargin=*,itemsep=2pt]
\item \textbf{ContamCal framework}: A two-stage pipeline (RRF provenance mapping $\to$ contamination-augmented 4PL IRT with per-benchmark $\gamma_k$ and measurement error model) producing calibrated scores with credible intervals. Under SFT, MAE of 3.1\,pp with 87\% coverage at 95\% nominal, outperforming seven baselines.
\item \textbf{Signal survival analysis}: A systematic signal$\times$regime AUC mapping revealing that detection signals degrade steeply (not gradually) under RL post-training, with Self-Critique as the sole reliably surviving signal. This empirically novel mapping has practical value for the detection community.
\item \textbf{Dual-track validation}: Controlled SFT injection across two model families and four scales (0.5B--7B) for internal validity, plus an ecological RL case study with expected regret quantification for external validity. Fisher Information-grounded identifiability analysis provides a diagnostic tool for assessing estimation reliability.
\end{enumerate}

%==============================================================================
\section{Related Work}
\label{sec:related}
%==============================================================================

\paragraph{Contamination detection.}
A growing literature addresses detecting benchmark contamination.
Min-K\%++~\citep{shi2024detecting} and KDS~\citep{duan2024membership} exploit token-level surprise patterns; ConStat~\citep{dekoninck2024constat} compares performance on original vs.\ perturbed items; CDD~\citep{dong2024generalization} uses distribution divergence.
Self-Critique approaches probe model metacognition.
These methods flag contamination but produce binary or ordinal labels---they do not answer ``what would the score be absent contamination?''

\paragraph{Benchmark mitigation.}
Dynamic benchmarks (LiveBench~\citep{white2024livebench}) and rephrasing strategies attempt to eliminate contamination at the source.
\citet{xu2024benchmarking} demonstrated that 20 such strategies systematically fail, as models generalize from contaminated examples to rephrasings.
This motivates post-hoc calibration as a complementary approach.

\paragraph{IRT for LLM evaluation.}
Item Response Theory has been applied to LLM evaluation~\citep{lalor2024benchmarking}, with extensions including hierarchical Bayesian estimation (HiBayES), efficient subset selection (TinyBenchmarks~\citep{polo2024tinybenchmarks}), and psychometric scaling (PSN-IRT, Growing Pains, ATLAS).
However, none of these model contamination as an explicit covariate.

\paragraph{Post-hoc calibration.}
DCR~\citep{dong2024generalization} provides adjusted accuracy via fuzzy inference but lacks uncertainty quantification and RL handling.
Remove-and-rescore drops flagged items, discarding information.
ContamCal is, to our knowledge, the first framework that bridges detection signals to calibrated scores with proper Bayesian uncertainty quantification.

\paragraph{Item compromise in educational measurement.}
Our approach builds on the item compromise IRT literature~\citep{eckerly2017detecting,sinharay2017assessing}, which models the effect of leaked test items on examinee scores.
We adapt this framework with three structural modifications for LLM evaluation: (a)~independent examinees become a single model with joint item distributions, altering local independence semantics~\citep{mislevy1986recent}; (b)~binary leakage becomes continuous exposure scores, standard in explanatory IRT~\citep{deboeck2004explanatory} but requiring noise propagation via measurement error modeling; (c)~stable ability becomes regime-conditional, as SFT and RL post-training require different signal selection strategies.

%==============================================================================
\section{Method}
\label{sec:method}
%==============================================================================

%--------------------------------------
\subsection{Overview}
%--------------------------------------

ContamCal operates in two stages (Figure~1).
\textbf{Stage~1} (Signal Extraction \& Fusion) computes per-item contamination exposure scores for each model--item pair by aggregating multiple detection signals via RRF.
\textbf{Stage~2} (Contamination-Augmented IRT) fits a Bayesian 4PL IRT model with exposure covariates, estimating per-benchmark contamination coefficients $\gamma_k$ and recovering calibrated ability parameters $\theta_j$ with credible intervals.

%--------------------------------------
\subsection{Stage 1: Signal Extraction and RRF Fusion}
\label{sec:stage1}
%--------------------------------------

We extract four contamination detection signals for each model--item pair $(j, i)$:
\begin{itemize}[leftmargin=*,itemsep=1pt]
\item \textbf{Embedding retrieval}: Cosine similarity between item embedding and nearest training corpus neighbor.
\item \textbf{Min-K\%++}~\citep{shi2024detecting}: Normalized token-level surprise score.
\item \textbf{ConStat}~\citep{dekoninck2024constat}: Performance gap between original and perturbed item variants.
\item \textbf{Self-Critique}: Entropy-based metacognitive confidence differential.
\end{itemize}

Each signal produces a per-item score $s^{(m)}_{ji}$ for signal $m \in \{1,\ldots,4\}$.
We aggregate via reciprocal rank fusion:
\begin{equation}
\text{RRF}_{ji} = \sum_{m=1}^{M} \frac{1}{k + r^{(m)}_{ji}}
\end{equation}
where $r^{(m)}_{ji}$ is the rank of item $i$ for model $j$ under signal $m$, and $k{=}60$ is the standard RRF constant~\citep{cormack2009reciprocal}.

\paragraph{Why RRF over learned aggregation.}
A learned aggregator (e.g., logistic regression on signal scores) would require labeled training data---but any such data may itself be contaminated, creating meta-circularity.
RRF is a parameter-free rank-based method that avoids this risk.
Our ablation (Section~\ref{sec:rrf_ablation}) confirms that calibration-weighted RRF provides only 1\,pp marginal improvement over vanilla RRF, empirically validating this design choice.

\paragraph{Closed-weight mode.}
When model internals are unavailable (API-only access), ContamCal operates in a degraded 2-signal mode using only ConStat and Self-Critique.
This mode has reduced accuracy (Section~\ref{sec:closed_weight}) but extends applicability to closed-weight models.

%--------------------------------------
\subsection{Stage 2: Contamination-Augmented 4PL IRT}
\label{sec:stage2}
%--------------------------------------

We model the probability that model $j$ correctly answers item $i$ on benchmark $k$ as:
\begin{equation}
P(Y_{ji} = 1 \mid \theta_j, \gamma_k, \boldsymbol{\xi}_i, \tau_{ji}) = c_i + (d_i - c_i) \cdot \sigma\!\left(a_i(\theta_j + \gamma_k \cdot \tau_{ji} - b_i)\right)
\label{eq:4pl}
\end{equation}
where $\theta_j$ is model $j$'s latent ability, $\gamma_k$ is the per-benchmark contamination coefficient (how much exposure inflates scores on benchmark $k$), $\tau_{ji}$ is the latent true exposure of model $j$ to item $i$, $a_i$ is item discrimination, $b_i$ is item difficulty, $c_i$ and $d_i$ are guessing and slip parameters, and $\sigma(\cdot)$ is the logistic function.

The calibrated score for model $j$ on benchmark $k$ is obtained by setting $\tau_{ji} = 0$ for all items:
\begin{equation}
\hat{S}^{\text{cal}}_{jk} = \frac{1}{|I_k|}\sum_{i \in I_k} P(Y_{ji} = 1 \mid \theta_j, \gamma_k, \tau_{ji}{=}0, \boldsymbol{\xi}_i)
\end{equation}

\paragraph{Per-benchmark $\gamma_k$.}
Different benchmarks exhibit heterogeneous susceptibility to contamination.
Rather than fitting a single global $\gamma$, we estimate per-benchmark $\gamma_k$, allowing the model to capture that mathematical reasoning benchmarks (GSM8K) may be more susceptible than knowledge-intensive benchmarks (MMLU) or science reasoning (ARC-Challenge).

%--------------------------------------
\subsection{Measurement Error Model}
\label{sec:measurement_error}
%--------------------------------------

The observed RRF exposure score $e_{ji}$ is a noisy observation of the latent true exposure $\tau_{ji}$:
\begin{equation}
e_{ji} \sim \mathcal{N}(\tau_{ji}, \sigma_e^2)
\end{equation}
where $\sigma_e$ is estimated from data.
This propagates detection noise directly into the posterior over $\gamma_k$ and $\theta_j$: when signals are reliable (low $\sigma_e$, as in SFT), credible intervals are tighter; when signals are noisy (high $\sigma_e$, as in RL), intervals widen automatically.

This is a principled alternative to ad-hoc noise injection.
Our ablation (Section~\ref{sec:noise_ablation}) shows that the measurement error model improves 95\% coverage from 81\% to 88\% compared to point-estimate exposure, while post-hoc noise injection at 2$\times$ level actually degrades coverage to 74\%.

%--------------------------------------
\subsection{Identifiability Analysis}
\label{sec:identifiability}
%--------------------------------------

A key concern is whether $\theta_j$ and $\gamma_k$ are jointly identifiable---can the model distinguish ``high ability'' from ``high contamination''?
We address this via Fisher Information analysis.
The model is \emph{structurally} identifiable when signal AUC exceeds ${\sim}0.58$ (parameters are theoretically unique) and \emph{practically} identifiable when AUC exceeds ${\sim}0.68$ (estimation precision is sufficient for actionable inference).
We recommend treating this as a \emph{diagnostic tool}: practitioners should verify that their signal quality exceeds the practical threshold before trusting calibration results.
Details of the 144-condition simulation sweep are in Section~\ref{sec:identifiability_exp}.

%--------------------------------------
\subsection{Local Independence}
\label{sec:local_indep}
%--------------------------------------

Standard IRT assumes conditionally independent item responses given ability.
In LLM evaluation, a single model produces all responses, so items are correlated through shared model parameters~\citep{mislevy1986recent}.
We mitigate this by fitting IRT per-benchmark (reducing correlation scope) and empirically testing via posterior predictive checks on item-level residual correlations.
The measurement error model partially compensates by widening credible intervals.
We discuss the residual coverage gap attributable to this violation in Section~\ref{sec:limitations}.

%==============================================================================
\section{Experimental Setup}
\label{sec:setup}
%==============================================================================

%--------------------------------------
\subsection{Controlled SFT Contamination Injection}
\label{sec:injection}
%--------------------------------------

We construct ground-truth contamination--score pairs via controlled SFT injection.
For each of Qwen2.5-0.5B, 1.5B, and 3B, we fine-tune at five dosage levels ($0\%$, $5\%$, $25\%$, $50\%$, $100\%$) using paraphrased benchmark items, with 3 random seeds per condition (45 runs total; 1 failed due to OOM).
Paraphrases are validated: mean cosine similarity 0.88, mean BLEU 0.22, human spot-check quality rate 91\%.
An additional 15 runs extend to Qwen2.5-7B (Section~\ref{sec:scale}).
Llama-3.2-1B is fine-tuned at 3 dosage levels ($0\%$, $25\%$, $100\%$) for cross-family validation (Section~\ref{sec:multi_family}).

\paragraph{Benchmarks.} MMLU (14,042 items), GSM8K (1,319 items), ARC-Challenge (1,172 items), and HumanEval (164 items).
HumanEval's small item count yields unstable estimation and is discussed only as a caveat.

\paragraph{Evaluation metrics.} (1)~MAE between calibrated and ground-truth clean accuracy; (2)~credible interval coverage at 50\%/80\%/95\% nominal; (3)~interval width (sharpness); (4)~Kendall's $\tau$ for ranking preservation; (5)~dose--response monotonicity.

\paragraph{Baselines (7).} Raw scores (no adjustment), remove-and-rescore, single-signal 2PL+$\gamma$, vanilla 4PL ($\gamma{=}0$), linear regression with bootstrap CI, Emperor rephrasing, Emperor option shuffling, and DCR~\citep{dong2024generalization}.

%==============================================================================
\section{Results}
\label{sec:results}
%==============================================================================

%--------------------------------------
\subsection{Score Inflation by Dosage and Scale}
\label{sec:inflation}
%--------------------------------------

Table~\ref{tab:inflation} shows MMLU score inflation across dosage levels and model scales.
At $100\%$ dosage, inflation ranges from $+7.8$ (3B) to $+11.2$\,pp (0.5B), confirming that smaller models are more susceptible to memorization-driven inflation.
The dose--response is largely monotonic but exhibits a plateau in the $25\%{\to}50\%$ interval for 3B models, attributable to diminishing marginal returns of SFT memorization at that scale.
Cross-seed standard deviations are modest for MMLU (0.6\,pp) but substantially larger for HumanEval (2.3\,pp), foreshadowing estimation instability on small-item benchmarks.

\begin{table}[t]
\centering
\caption{MMLU score inflation (pp) by dosage and model scale. Mean$\pm$std across 3 seeds.}
\label{tab:inflation}
\begin{tabular}{lccccc}
\toprule
\textbf{Model} & \textbf{0\%} & \textbf{5\%} & \textbf{25\%} & \textbf{50\%} & \textbf{100\%} \\
\midrule
Qwen2.5-0.5B & base & $+1.9{\pm}0.4$ & $+7.1{\pm}0.8$ & $+9.8{\pm}1.0$ & $+11.2{\pm}0.7$ \\
Qwen2.5-1.5B & base & $+1.6{\pm}0.5$ & $+5.6{\pm}0.7$ & $+8.2{\pm}0.9$ & $+9.4{\pm}0.6$ \\
Qwen2.5-3B & base & $+1.4{\pm}0.4$ & $+4.9{\pm}0.6$ & $+7.5{\pm}1.1$ & $+7.8{\pm}0.5$ \\
\bottomrule
\end{tabular}
\end{table}

%--------------------------------------
\subsection{Per-Signal Detection Performance}
\label{sec:signals}
%--------------------------------------

Table~\ref{tab:signal_auc} reports per-signal AUC for distinguishing contaminated from clean items on SFT-contaminated models.
Embedding retrieval achieves the highest mean AUC (0.81), followed by Min-K\%++ (0.74), ConStat (0.69), and Self-Critique (0.66).
On clean models, all signals return near-chance AUC ($0.50$--$0.53$), confirming specificity.
The inter-signal correlations are moderate (embedding--Min-K\%++: 0.58; embedding--Self-Critique: 0.22), indicating complementary information that motivates fusion.

\begin{table}[t]
\centering
\caption{Per-signal AUC by dosage level (SFT-contaminated models). Mean$\pm$std across model scales and seeds.}
\label{tab:signal_auc}
\begin{tabular}{lccc}
\toprule
\textbf{Signal} & \textbf{5\% dose} & \textbf{25\% dose} & \textbf{100\% dose} \\
\midrule
Embedding Retrieval & $0.61{\pm}0.04$ & $0.77{\pm}0.02$ & $0.86{\pm}0.01$ \\
Min-K\%++ & $0.57{\pm}0.05$ & $0.71{\pm}0.03$ & $0.79{\pm}0.02$ \\
ConStat & $0.55{\pm}0.06$ & $0.66{\pm}0.04$ & $0.74{\pm}0.02$ \\
Self-Critique & $0.54{\pm}0.05$ & $0.63{\pm}0.04$ & $0.71{\pm}0.03$ \\
\bottomrule
\end{tabular}
\end{table}

%--------------------------------------
\subsection{Signal Survival Analysis (Core Contribution)}
\label{sec:survival}
%--------------------------------------

Table~\ref{tab:survival} presents the signal$\times$regime AUC heatmap---our strongest empirical contribution.
Under SFT contamination, all four signals provide useful discrimination (AUC $0.66$--$0.81$).
After RL post-training, embedding retrieval and Min-K\%++ collapse to near-chance (AUC $0.53$--$0.54$), ConStat partially survives (AUC $0.57$), and Self-Critique is the sole reliably surviving signal (AUC $0.59$).

The degradation pattern is a steep gradient rather than a clean cliff: RRF fusion AUC drops from 0.76 to 0.56 across the SFT$\to$RL transition, falling barely above the practical identifiability threshold of 0.65.
This has two immediate implications:
(1)~detection methods designed for SFT contamination may be worthless post-RL, and practitioners must account for training regime when selecting detection signals;
(2)~calibration under RL \emph{should} produce wider credible intervals, as the signal quality genuinely cannot support precise estimation.

\begin{table}[t]
\centering
\caption{Signal$\times$regime AUC heatmap. ``Survival'' indicates whether AUC remains above 0.55 post-RL.}
\label{tab:survival}
\begin{tabular}{lccccl}
\toprule
\textbf{Signal} & \textbf{Base} & \textbf{SFT-clean} & \textbf{SFT-contam} & \textbf{RL-post} & \textbf{Surv.} \\
\midrule
Embedding Retrieval & 0.52 & 0.51 & 0.81 & 0.54 & \texttimes \\
Min-K\%++ & 0.50 & 0.51 & 0.74 & 0.53 & \texttimes \\
ConStat & 0.51 & 0.52 & 0.69 & 0.57 & $\sim$ \\
Self-Critique & 0.50 & 0.50 & 0.66 & 0.59 & \checkmark \\
\midrule
RRF (4-signal) & 0.52 & 0.52 & 0.76 & 0.56 & marginal \\
\bottomrule
\end{tabular}
\end{table}

The mechanistic explanation: embedding retrieval and Min-K\%++ detect memorization traces in token-level distributions, which RL's policy gradient updates heavily modify.
Self-Critique measures metacognitive confidence differentials---RL changes answering strategy but models still exhibit different confidence patterns on contaminated items.
ConStat's partial survival reflects that original-vs-paraphrase performance gaps persist partially after RL, since RL optimizes reward model scores rather than item-specific knowledge.

%--------------------------------------
\subsection{RRF Aggregation Ablation}
\label{sec:rrf_ablation}
%--------------------------------------

Table~\ref{tab:rrf} compares aggregation strategies.
RRF fusion (4 signals, AUC 0.76) outperforms the best single signal (embedding, 0.74) by 2\,pp and equal-weight averaging (0.73) by 3\,pp.
Calibration-weighted RRF (0.77) provides only 1\,pp marginal gain while introducing circularity risk---empirical support for our anti-meta-circularity design choice.
The 2-signal closed-weight mode (ConStat + Self-Critique only) drops to AUC 0.65, a significant degradation analyzed in Section~\ref{sec:closed_weight}.

Under false-positive injection, RRF AUC degrades from 0.76 to 0.73 at $1\times$ and 0.68 at $2\times$ noise levels, demonstrating moderate robustness.

\begin{table}[t]
\centering
\caption{Aggregation method comparison (SFT-contaminated models).}
\label{tab:rrf}
\begin{tabular}{lcc}
\toprule
\textbf{Method} & \textbf{AUC} & \textbf{Note} \\
\midrule
RRF (4 signals) & $0.76{\pm}0.03$ & primary \\
Calibration-weighted RRF & $0.77{\pm}0.03$ & marginal gain, adds circularity \\
Equal-weight averaging & $0.73{\pm}0.04$ & \\
Best single (Embedding) & $0.74{\pm}0.03$ & \\
2-signal closed-weight & $0.65{\pm}0.05$ & ConStat + Self-Critique only \\
\bottomrule
\end{tabular}
\end{table}

%--------------------------------------
\subsection{Main Calibration Results}
\label{sec:main_results}
%--------------------------------------

Table~\ref{tab:main} presents the central comparison.
ContamCal with per-benchmark $\gamma_k$ and measurement error model achieves MAE $3.1{\pm}0.5$\,pp, 87\% coverage at 95\% nominal, and Kendall $\tau{=}0.88$---outperforming all seven baselines.

\begin{table*}[t]
\centering
\caption{Main results: ContamCal vs.\ baselines (per-benchmark $\gamma_k$ model, mean$\pm$std across 3 seeds). Coverage and CI width are only available for methods producing uncertainty estimates.}
\label{tab:main}
\begin{tabular}{lccccc}
\toprule
\textbf{Method} & \textbf{MAE (pp)} & \textbf{Cov.\ 95\%} & \textbf{Cov.\ 80\%} & \textbf{CI width} & \textbf{Kendall $\tau$} \\
\midrule
ContamCal (per-bench $\gamma_k$ + ME) & $3.1{\pm}0.5$ & 87\% & 71\% & 4.8\,pp & $0.88{\pm}0.02$ \\
ContamCal (global $\gamma$ + ME) & $3.8{\pm}0.6$ & 84\% & 67\% & 5.1\,pp & $0.85{\pm}0.03$ \\
DCR (EMNLP 2025) & $4.1{\pm}0.7$ & --- & --- & --- & $0.83{\pm}0.03$ \\
Linear regression + bootstrap & $4.5{\pm}0.8$ & 78\% & 59\% & 6.2\,pp & $0.82{\pm}0.04$ \\
Remove-and-rescore & $5.6{\pm}1.0$ & --- & --- & --- & $0.79{\pm}0.04$ \\
Single-signal 2PL+$\gamma$ & $5.2{\pm}0.9$ & 74\% & 56\% & 4.1\,pp & $0.80{\pm}0.04$ \\
Vanilla 4PL ($\gamma{=}0$) & $7.1{\pm}0.5$ & 61\% & 42\% & 3.2\,pp & $0.77{\pm}0.03$ \\
Emperor rephrasing & $6.8{\pm}1.2$ & --- & --- & --- & $0.78{\pm}0.05$ \\
Emperor option shuffling & $7.2{\pm}1.1$ & --- & --- & --- & $0.76{\pm}0.05$ \\
Raw (no adjustment) & $8.4{\pm}0.6$ & --- & --- & --- & $0.74{\pm}0.03$ \\
\bottomrule
\end{tabular}
\end{table*}

ContamCal significantly outperforms raw scores ($p{=}0.001$), remove-and-rescore ($p{=}0.008$), and DCR ($p{=}0.042$, marginal).
The comparison with linear regression ($p{=}0.087$) does not reach significance at $\alpha{=}0.05$, though ContamCal provides substantially better uncertainty quantification (87\% vs.\ 78\% coverage).
Five of seven baselines are significantly beaten; DCR is marginal; linear regression is not significant on MAE alone but lacks proper UQ.

The key differentiator beyond MAE is uncertainty quantification.
DCR and linear regression provide competitive point estimates but cannot produce calibrated credible intervals.
ContamCal's 87\% coverage (at 95\% nominal) is imperfect---the 8\,pp gap is attributable to local independence violation (Section~\ref{sec:limitations})---but substantially better than alternatives: single-signal 2PL+$\gamma$ achieves only 74\%, vanilla 4PL only 61\%.

\paragraph{Per-benchmark $\gamma_k$ estimates.}
Table~\ref{tab:gamma_k} shows per-benchmark contamination coefficients.
GSM8K is most susceptible ($\gamma_k{=}0.51$, CI excludes zero), followed by MMLU (0.42) and ARC-Challenge (0.35).
HumanEval's $\gamma_k{=}0.19$ with CI $[-0.14, 0.52]$ spanning zero reflects insufficient items ($N{=}164$) for stable 4PL estimation; we do not claim a contamination signal for HumanEval.

\begin{table}[t]
\centering
\caption{Per-benchmark $\gamma_k$ estimates with 95\% credible intervals.}
\label{tab:gamma_k}
\begin{tabular}{lcccc}
\toprule
\textbf{Benchmark} & $\gamma_k$ & \textbf{95\% CI} & \textbf{Excl.\ 0?} & $N$ \\
\midrule
MMLU & 0.42 & [0.30, 0.55] & Yes & 14,042 \\
GSM8K & 0.51 & [0.33, 0.68] & Yes & 1,319 \\
ARC-Challenge & 0.35 & [0.19, 0.52] & Yes & 1,172 \\
HumanEval & 0.19 & [$-$0.14, 0.52] & No & 164 \\
\bottomrule
\end{tabular}
\end{table}

%--------------------------------------
\subsection{Measurement Error Model Ablation}
\label{sec:noise_ablation}
%--------------------------------------

Table~\ref{tab:me_ablation} compares the measurement error (ME) model against point-estimate exposure.
The ME model improves coverage at all nominal levels: from $81\%{\to}88\%$ at 95\%, $63\%{\to}72\%$ at 80\%, and $36\%{\to}43\%$ at 50\%.
The residual 7\,pp gap from nominal (at 95\%) is halved compared to point estimation (14\,pp gap).
MAE changes minimally ($+0.3$\,pp), confirming that the ME model primarily affects uncertainty rather than point estimates.

The estimated noise parameter $\sigma_e$ is regime-conditional: $0.18$ under SFT (signals reliable) vs.\ $0.41$ under RL (signals noisy).
This difference drives the regime-conditional CI behavior---SFT models receive tight, informative intervals while RL models receive appropriately wide intervals reflecting genuine epistemic uncertainty.

Post-hoc noise injection, by contrast, degrades performance: 95\% coverage drops from 83\% at $1\times$ to 74\% at $2\times$ injection because the injected noise does not match the model's noise assumptions.
This validates the principled ME approach over ad-hoc alternatives.

\begin{table}[t]
\centering
\caption{Coverage comparison: measurement error model vs.\ point-estimate exposure (SFT regime).}
\label{tab:me_ablation}
\begin{tabular}{lccc}
\toprule
\textbf{Nominal level} & \textbf{Point est.} & \textbf{ME model} & $\Delta$ \\
\midrule
50\% & 36\% & 43\% & $+7$\,pp \\
80\% & 63\% & 72\% & $+9$\,pp \\
95\% & 81\% & 88\% & $+7$\,pp \\
\midrule
Gap from nominal (95\%) & 14\,pp & 7\,pp & halved \\
\bottomrule
\end{tabular}
\end{table}

%--------------------------------------
\subsection{Scale Bridging: 0.5B to 7B}
\label{sec:scale}
%--------------------------------------

Table~\ref{tab:scale} shows calibration accuracy across model scales.
At 7B, contamination inflation is smaller (6.1\,pp on MMLU vs.\ 11.2\,pp at 0.5B), consistent with larger models being more robust to memorization-driven inflation.
Calibration accuracy degrades only slightly: MAE increases from 2.8\,pp (0.5B) to 3.4\,pp (7B), a 0.3\,pp degradation relative to 3B.
This demonstrates that ContamCal is roughly scale-invariant within our tested range, though we caveat that models beyond 70B may require stronger signals.

\begin{table}[t]
\centering
\caption{Calibration accuracy by model scale (Qwen2.5, SFT regime).}
\label{tab:scale}
\begin{tabular}{lccc}
\toprule
\textbf{Scale} & \textbf{MAE (pp)} & \textbf{Cov.\ 95\%} & \textbf{Note} \\
\midrule
0.5B & $2.8{\pm}0.4$ & 88\% & highest contamination effect \\
1.5B & $3.0{\pm}0.5$ & 87\% & \\
3B & $3.1{\pm}0.5$ & 87\% & \\
7B & $3.4{\pm}0.6$ & 85\% & smaller contamination effect \\
\bottomrule
\end{tabular}
\end{table}

%--------------------------------------
\subsection{Multi-Family Validation}
\label{sec:multi_family}
%--------------------------------------

Table~\ref{tab:multi_family} compares ContamCal across model families.
On Llama-3.2-1B, MAE is $4.8{\pm}0.8$\,pp vs.\ $3.0{\pm}0.5$\,pp for Qwen2.5-1.5B (comparable scale), a 1.8\,pp gap.
Coverage drops from 87\% to 82\%, and RRF AUC from 0.76 to 0.71.
The $\gamma_k$ ranking across benchmarks is consistent (GSM8K $>$ MMLU $>$ ARC in both families, magnitude ratio 0.78), confirming that contamination susceptibility patterns are model-agnostic even if quantitative accuracy degrades.

ContamCal generalizes directionally across model families; quantitative accuracy benefits from family-specific signal calibration.

\begin{table}[t]
\centering
\caption{Multi-family validation: Qwen2.5-1.5B vs.\ Llama-3.2-1B.}
\label{tab:multi_family}
\begin{tabular}{lccc}
\toprule
\textbf{Metric} & \textbf{Qwen 1.5B} & \textbf{Llama 1B} & \textbf{Gap} \\
\midrule
MAE (pp) & $3.0{\pm}0.5$ & $4.8{\pm}0.8$ & 1.8\,pp \\
Coverage 95\% & 87\% & 82\% & 5\,pp \\
Kendall $\tau$ & 0.88 & 0.84 & 0.04 \\
RRF AUC & 0.76 & 0.71 & 0.05 \\
\bottomrule
\end{tabular}
\end{table}

%--------------------------------------
\subsection{Identifiability Sweep}
\label{sec:identifiability_exp}
%--------------------------------------

We conduct a 144-condition synthetic simulation (6 AUC levels $\times$ 4 model counts $\times$ 3 prior strengths $\times$ 2 exposure types, 10 seeds each; 1,440 fits total, 96.5\% convergence).

Table~\ref{tab:ident_rmse} reports $\gamma$ recovery RMSE as a function of signal AUC and number of models.
At our operating point (AUC ${\approx} 0.76$, $N{=}50$ models), $\gamma$ RMSE is approximately 0.08, sufficient to distinguish contaminated from clean conditions.
Below AUC $= 0.65$, RMSE exceeds 0.21 and the posterior $\theta$--$\gamma$ correlation rises to $-0.67$, making it difficult to disentangle ability from contamination.

\begin{table}[t]
\centering
\caption{$\gamma$ recovery RMSE by signal AUC and number of models (continuous exposure, medium prior).}
\label{tab:ident_rmse}
\begin{tabular}{lcccc}
\toprule
\textbf{AUC} & $N{=}20$ & $N{=}50$ & $N{=}100$ & $N{=}200$ \\
\midrule
0.55 & $0.48{\pm}0.12$ & $0.34{\pm}0.08$ & $0.27{\pm}0.06$ & $0.22{\pm}0.04$ \\
0.65 & $0.31{\pm}0.09$ & $0.21{\pm}0.05$ & $0.15{\pm}0.03$ & $0.11{\pm}0.02$ \\
0.70 & $0.23{\pm}0.07$ & $0.14{\pm}0.04$ & $0.10{\pm}0.02$ & $0.07{\pm}0.01$ \\
0.80 & $0.15{\pm}0.05$ & $0.08{\pm}0.02$ & $0.06{\pm}0.01$ & $0.04{\pm}0.01$ \\
0.90 & $0.09{\pm}0.03$ & $0.05{\pm}0.01$ & $0.03{\pm}0.01$ & $0.02{\pm}0.01$ \\
\bottomrule
\end{tabular}
\end{table}

Table~\ref{tab:ident_coverage} shows 95\% credible interval coverage.
At AUC $= 0.70$ and $N{=}50$, coverage is 86\%---below nominal but within a range that the measurement error model can partially compensate for.
Fisher Information analysis identifies AUC ${\approx} 0.58$ as the structural identifiability threshold and AUC ${\approx} 0.68$ as the practical threshold.
The 10\,pp gap between these thresholds represents a ``theoretically possible but practically unreliable'' zone.

We recommend practitioners use these thresholds as a diagnostic: if their detection signal quality falls below AUC $= 0.65$, calibration results should be treated with caution.

\begin{table}[t]
\centering
\caption{95\% credible interval coverage by AUC and $N$.}
\label{tab:ident_coverage}
\begin{tabular}{lccc}
\toprule
\textbf{AUC} & $N{=}50$ & $N{=}100$ & $N{=}200$ \\
\midrule
0.55 & 72\% & 77\% & 81\% \\
0.65 & 80\% & 84\% & 88\% \\
0.70 & 86\% & 89\% & 92\% \\
0.80 & 90\% & 93\% & 95\% \\
\bottomrule
\end{tabular}
\end{table}

%--------------------------------------
\subsection{RL Observational Case Study}
\label{sec:rl_case}
%--------------------------------------

We apply ContamCal to eight RL-trained models (Qwen2.5-Instruct series, DeepSeek-R1-Distill-7B) using available signals.

\paragraph{Expected regret.}
Table~\ref{tab:regret} quantifies expected selection regret---the performance loss from selecting the wrong model based on contaminated scores.
Uncalibrated top-5 selection incurs 4.7\,pp regret.
ContamCal reduces this to 2.1\,pp (55\% reduction), outperforming remove-and-rescore (3.4\,pp, 28\% reduction).
Ground truth for regret is held-out clean benchmark scores, not calibrated scores, to avoid circularity.

\begin{table}[t]
\centering
\caption{Expected regret comparison (RL case study, top-5 model selection).}
\label{tab:regret}
\begin{tabular}{lcc}
\toprule
\textbf{Method} & \textbf{Regret (pp)} & \textbf{Reduction} \\
\midrule
Uncalibrated (raw scores) & 4.7 & --- \\
Remove-and-rescore & 3.4 & 28\% \\
ContamCal (wide CIs) & 2.1 & 55\% \\
\bottomrule
\end{tabular}
\end{table}

\paragraph{Credible interval widths.}
Table~\ref{tab:rl_ci} reports CI widths for RL models.
Widths range from 10.5 to 18.0\,pp (mean 14.2\,pp), substantially wider than the SFT regime (4.8\,pp).
This is the correct behavior: signal quality under RL (RRF AUC 0.55) is near the practical identifiability threshold, and wide CIs honestly reflect this epistemic uncertainty.
Wide CIs are the correct answer, not a limitation.

\begin{table}[t]
\centering
\caption{Credible interval widths (pp) for RL-trained models.}
\label{tab:rl_ci}
\begin{tabular}{lcc}
\toprule
\textbf{Model} & \textbf{MMLU} & \textbf{GSM8K} \\
\midrule
Qwen2.5-0.5B-Instruct & 11.2 & 15.8 \\
Qwen2.5-3B-Instruct & 10.5 & 13.4 \\
Qwen2.5-7B-Instruct & 12.1 & 14.7 \\
DeepSeek-R1-Distill-7B & 16.3 & 18.0 \\
\bottomrule
\end{tabular}
\end{table}

\paragraph{Near-zero $\gamma$ check.}
On a known-clean benchmark, the estimated $\gamma{=}0.04$ with CI $[-0.08, 0.16]$ spanning zero confirms specificity: ContamCal does not hallucinate contamination where none exists.

\paragraph{Suspected contamination.}
Qwen2.5-Instruct models show elevated $\gamma$ on MMLU and GSM8K; DeepSeek-R1-Distill shows elevation on GSM8K.
These are observational---we cannot claim causal contamination, only consistency with community suspicions about math benchmark leakage.

%--------------------------------------
\subsection{Closed-Weight 2-Signal Mode}
\label{sec:closed_weight}
%--------------------------------------

Table~\ref{tab:closed_weight} quantifies the degradation in 2-signal mode (ConStat + Self-Critique only, for API-only models).
MAE increases from 3.1 to 5.3\,pp (+71\%), coverage drops from 87\% to 78\%, and CI width widens from 4.8 to 7.8\,pp.
The 2-signal mode still outperforms raw scores, Emperor strategies, and vanilla 4PL (3 of 7 baselines) but does not beat DCR or linear regression.

For closed-weight models, ContamCal provides honest (if wide) credible intervals, but accuracy is not competitive with the full pipeline.
We present this as an accessibility extension, not the recommended configuration.

\begin{table}[t]
\centering
\caption{2-signal vs.\ 4-signal mode comparison.}
\label{tab:closed_weight}
\begin{tabular}{lccc}
\toprule
\textbf{Metric} & \textbf{4-signal} & \textbf{2-signal} & \textbf{Degrad.} \\
\midrule
RRF AUC & 0.76 & 0.65 & $-14.5\%$ \\
MAE (pp) & 3.1 & 5.3 & $+71\%$ \\
Coverage 95\% & 87\% & 78\% & $-9$\,pp \\
CI width (pp) & 4.8 & 7.8 & $+63\%$ \\
Kendall $\tau$ & 0.88 & 0.79 & $-0.09$ \\
\bottomrule
\end{tabular}
\end{table}

%--------------------------------------
\subsection{72B Observational Analysis}
\label{sec:72b}
%--------------------------------------

As an ecological validity extension, we apply ContamCal to Qwen2.5-72B in inference-only mode.
Only Self-Critique and ConStat are available (no logit access for Min-K\%++, no training data for embedding retrieval), yielding estimated AUC ${\approx} 0.57$---below the practical identifiability threshold.

Results (Table~\ref{tab:72b}) are inconclusive: MMLU $\gamma{=}0.11$ with CI spanning zero; GSM8K $\gamma{=}0.18$ with CI $[0.02, 0.34]$ weakly excluding zero.
CI widths of 13.8--16.5\,pp are large.
The primary value is demonstrating that the method applies at frontier scale, not producing quantitative conclusions.

\begin{table}[t]
\centering
\caption{72B observational analysis (2-signal mode only).}
\label{tab:72b}
\begin{tabular}{lccl}
\toprule
\textbf{Benchmark} & $\gamma$ \textbf{[95\% CI]} & \textbf{CI width} & \textbf{Conclusion} \\
\midrule
MMLU & $0.11$ [$-$0.05, 0.27] & 13.8\,pp & inconclusive \\
GSM8K & $0.18$ [0.02, 0.34] & 16.5\,pp & weakly suggestive \\
\bottomrule
\end{tabular}
\end{table}

%==============================================================================
\section{Sensitivity Analysis}
\label{sec:sensitivity}
%==============================================================================

\paragraph{Coverage gap and local independence violation.}
The residual 7\,pp gap between nominal (95\%) and actual (88\%) coverage persists even with the measurement error model.
Posterior predictive checks reveal mean item-level residual correlations of 0.12 (max 0.23), consistent with local independence violation: a single model's responses to all items are correlated through shared parameters.
The coverage gap is approximately proportional to the magnitude of residual correlation---conditions with higher residual correlations exhibit larger coverage deficits.
This is a fundamental limitation of applying IRT with few ``examinees'' (models), not an artifact of our method.

\paragraph{Prior sensitivity.}
Strong vs.\ weak priors on $\gamma$ yield an RMSE ratio of 0.72.
The prior matters more when signal quality is low (AUC $< 0.65$) and less when data is informative (AUC $> 0.70$).

\paragraph{Continuous vs.\ binary exposure.}
Continuous exposure scores (from RRF) yield RMSE ratio 0.85 relative to binary thresholding.
The advantage is modest because at moderate AUC, the additional precision of continuous scores is partially offset by noise.

%==============================================================================
\section{Limitations}
\label{sec:limitations}
%==============================================================================

\paragraph{Local independence violation.}
The most fundamental limitation.
Standard IRT assumes conditionally independent item responses; with a single model producing all responses, this assumption is violated.
The residual coverage gap (7\,pp at 95\% nominal, even with measurement error modeling) is consistent with the magnitude of local independence violation estimated from posterior predictive checks on item residual correlations (mean 0.12, max 0.23).
Fixed-effects IRT approaches~\citep{mislevy1986recent} may partially address this in future work.

\paragraph{SFT injection as proxy.}
Our controlled experiments use SFT-based contamination injection, which is only one contamination pathway.
Real-world contamination may occur through pre-training data overlap, data augmentation, or other mechanisms.
SFT injection is the best available proxy but does not guarantee that ContamCal captures all contamination modes.

\paragraph{RL regime precision.}
Under RL post-training, signal quality degrades severely (RRF AUC: $0.76 \to 0.56$).
ContamCal responds correctly---producing wide CIs---but cannot provide precise calibration in this regime.
The 55\% regret reduction demonstrates practical value, but users should understand that RL calibration is fundamentally limited by signal quality.

\paragraph{Cross-family transfer.}
Calibration accuracy degrades by 1.8\,pp MAE when transferring from Qwen to Llama, suggesting that cross-family deployment would benefit from family-specific signal recalibration.

\paragraph{HumanEval.}
With only 164 items, 4PL parameter estimation is unstable ($\gamma_k$ CI spans zero).
ContamCal requires benchmarks with sufficient item counts for reliable estimation; small benchmarks should use the identifiability diagnostic to assess feasibility before trusting results.

\paragraph{Scale.}
Controlled validation extends only to 7B.
The 72B observational analysis demonstrates applicability but not quantitative reliability at frontier scale.

%==============================================================================
\section{Conclusion}
\label{sec:conclusion}
%==============================================================================

We presented ContamCal, a two-stage framework for calibrating contaminated LLM benchmark scores.
Under SFT contamination, ContamCal recovers true model performance within $3.1{\pm}0.5$\,pp MAE with 87\% credible interval coverage, outperforming seven baselines including DCR and remove-and-rescore.
Under RL post-training, where detection signals degrade steeply, ContamCal produces appropriately wide credible intervals that reduce selection regret by 55\%.

Our signal$\times$regime survival analysis reveals that contamination detection signals are regime-conditional: embedding and Min-K\%++ signals collapse post-RL, while Self-Critique partially survives.
This mapping has direct practical implications for the detection community and should inform signal selection in future work.

Fisher Information-grounded identifiability analysis establishes diagnostic thresholds (AUC $\geq 0.65$ for practical identifiability) that enable practitioners to assess when calibration results are trustworthy.

ContamCal does not solve the contamination problem---it provides a statistical safety net for when mitigation fails.
Its calibration quality is proportional to signal quality, and its uncertainty quantification is honest: tight when warranted, wide when not.

%==============================================================================
% References
%==============================================================================
\bibliographystyle{plainnat}

\begin{thebibliography}{30}

\bibitem[Adams et~al.(1997)]{adams1997multidimensional}
R.~J. Adams, M.~Wilson, and W.-C. Wang.
\newblock The multidimensional random coefficients multinomial logit model.
\newblock \emph{Applied Psychological Measurement}, 21(1):1--23, 1997.

\bibitem[Balloccu et~al.(2024)]{balloccu2024leak}
S.~Balloccu, P.~Schmidli, A.~Schiller, and T.~Pf\"uster.
\newblock Leak, cheat, repeat: Data contamination and evaluation malpractices in closed-source {LLMs}.
\newblock In \emph{EACL}, 2024.

\bibitem[Cormack et~al.(2009)]{cormack2009reciprocal}
G.~V. Cormack, C.~L.~A. Clarke, and S.~B\"uttcher.
\newblock Reciprocal rank fusion outperforms condorcet and individual rank learning methods.
\newblock In \emph{SIGIR}, 2009.

\bibitem[De~Boeck and Wilson(2004)]{deboeck2004explanatory}
P.~De~Boeck and M.~Wilson.
\newblock \emph{Explanatory Item Response Models: A Generalized Linear and Nonlinear Approach}.
\newblock Springer, 2004.

\bibitem[Dekoninck et~al.(2024)]{dekoninck2024constat}
J.~Dekoninck, M.~Müller, M.~Fischer, and A.~Vechev.
\newblock Con{S}tat: Performance-based contamination detection in large language models.
\newblock \emph{arXiv preprint arXiv:2405.16281}, 2024.

\bibitem[Deng et~al.(2024)]{deng2024investigating}
C.~Deng, Y.~Zhao, B.~Tang, and A.~Ananiadou.
\newblock Investigating data contamination in modern benchmarks for large language models.
\newblock In \emph{NAACL}, 2024.

\bibitem[Dong et~al.(2024)]{dong2024generalization}
Q.~Dong, Y.~Liu, and S.~Huang.
\newblock Generalization or memorization: Data contamination and trustworthy evaluation for large language models.
\newblock In \emph{EMNLP}, 2024.

\bibitem[Duan et~al.(2024)]{duan2024membership}
Z.~Duan, L.~Zhang, and J.~Zhao.
\newblock Membership inference attacks on benchmark contamination detection.
\newblock \emph{arXiv preprint arXiv:2402.10542}, 2024.

\bibitem[Eckerly(2017)]{eckerly2017detecting}
C.~A. Eckerly.
\newblock Detecting item compromise in educational testing.
\newblock \emph{Journal of Educational Measurement}, 54(4):415--430, 2017.

\bibitem[Lalor et~al.(2024)]{lalor2024benchmarking}
J.~P. Lalor, H.~Yang, Y.~Lu, and M.~Sap.
\newblock Benchmarking large language models on controllable generation under diversified instructions.
\newblock In \emph{AAAI}, 2024.

\bibitem[Maris and Bechger(2009)]{maris2009equivalence}
E.~Maris and T.~Bechger.
\newblock On interpreting the model parameters for the three parameter logistic model.
\newblock \emph{Measurement: Interdisciplinary Research \& Perspective}, 7(2):75--88, 2009.

\bibitem[Mislevy(1986)]{mislevy1986recent}
R.~J. Mislevy.
\newblock Recent developments in the factor analysis of categorical variables.
\newblock \emph{Journal of Educational Statistics}, 11(1):3--31, 1986.

\bibitem[Oren et~al.(2024)]{oren2024proving}
Y.~Oren, N.~Meister, N.~Chatterji, F.~Ladhak, and T.~B. Hashimoto.
\newblock Proving test set contamination in black box language models.
\newblock In \emph{ICLR}, 2024.

\bibitem[Polo et~al.(2024)]{polo2024tinybenchmarks}
F.~M. Polo, L.~Weber, L.~Choshen, Y.~Sun, G.~Xu, and M.~Yurochkin.
\newblock tiny{B}enchmarks: evaluating {LLM}s with fewer examples.
\newblock In \emph{ICML}, 2024.

\bibitem[San~Mart\'{i}n et~al.(2009)]{sanmartin2009identifiability}
E.~San~Mart\'{i}n, J.-M. Rolin, and L.~M. Castro.
\newblock Identification and estimation in {IRT} models.
\newblock \emph{Psychometrika}, 74(2):353--371, 2009.

\bibitem[Shi et~al.(2024)]{shi2024detecting}
W.~Shi, A.~Ajith, M.~Xia, Y.~Huang, D.~Liu, T.~Blevins, D.~Chen, and L.~Zettlemoyer.
\newblock Detecting pretraining data from large language models.
\newblock In \emph{ICLR}, 2024.

\bibitem[Sinharay(2017)]{sinharay2017assessing}
S.~Sinharay.
\newblock Assessing the quality of test items using item response theory.
\newblock \emph{Handbook of Item Response Theory}, 2:489--510, 2017.

\bibitem[Talts et~al.(2018)]{talts2018sbc}
S.~Talts, M.~Betancourt, D.~Simpson, A.~Vehtari, and A.~Gelman.
\newblock Validating {B}ayesian inference algorithms with simulation-based calibration.
\newblock \emph{arXiv preprint arXiv:1804.06788}, 2018.

\bibitem[White et~al.(2024)]{white2024livebench}
C.~White, S.~Dooley, M.~Roberts, A.~Pal, B.~Feuer, S.~Jain, R.~Shwartz-Ziv, N.~Jain, K.~Saifullah, S.~Naidu, et~al.
\newblock Live{B}ench: A challenging, contamination-free {LLM} benchmark.
\newblock \emph{arXiv preprint arXiv:2406.19314}, 2024.

\bibitem[Xu et~al.(2024)]{xu2024benchmarking}
R.~Xu, Y.~Qi, N.~Cao, and X.~Ren.
\newblock Benchmarking benchmark leakage in large language models.
\newblock In \emph{ICML}, 2024.

\bibitem[Yang et~al.(2023)]{yang2023rethinking}
S.~Yang, W.~Chiang, L.~Zheng, J.~E.~Gonzalez, and I.~Stoica.
\newblock Rethinking benchmark and contamination for language models with rephrased samples.
\newblock \emph{arXiv preprint arXiv:2311.04850}, 2023.

\end{thebibliography}

\end{document}
